\section{Dominant morphisms}
\label{section-dominant}

\noindent
The definition of a morphism of schemes being dominant is a little
different from what you might expect if you are used to the notion
of a dominant morphism of varieties.

\begin{definition}
\label{definition-dominant}
A morphism $f : X \to S$ of schemes is called {\it dominant} if the
image of $f$ is a dense subset of $S$.
\end{definition}

\noindent
So for example, if $k$ is an infinite field and $\lambda_1, \lambda_2, \ldots$
is a countable collection of distinct elements of $k$, then the morphism
$$
\coprod\nolimits_{i = 1, 2, \ldots } \Spec(k)
\longrightarrow
\Spec(k[x])
$$
with $i$th factor mapping to the point $x = \lambda_i$ is dominant.

\begin{lemma}
\label{lemma-generic-points-in-image-dominant}
Let $f : X \to S$ be a morphism of schemes.
If every generic point of every irreducible component of $S$
is in the image of $f$, then $f$ is dominant.
\end{lemma}

\begin{proof}
This is a topological fact which follows directly from the fact that
the topological space underlying a scheme is sober, see
Schemes, Lemma \ref{schemes-lemma-scheme-sober}, and that
every point of $S$ is contained in an irreducible component of
$S$, see Topology, Lemma \ref{topology-lemma-irreducible}.
\end{proof}

\noindent
The expectation that morphisms are dominant only if generic points of the
target are in the image does hold if the morphism is quasi-compact.

\begin{lemma}
\label{lemma-quasi-compact-dominant}
\begin{slogan}
Morphisms whose image contains the generic points are dominant
\end{slogan}
Let $f : X \to S$ be a quasi-compact morphism of schemes.
Then $f$ is dominant if and only if for every irreducible
component $Z \subset S$ the generic point of $Z$ is in the
image of $f$.
\end{lemma}

\begin{proof}
Let $V \subset S$ be an affine open.
Because $f$ is quasi-compact we may choose finitely many affine
opens $U_i \subset f^{-1}(V)$, $i = 1, \ldots, n$ covering
$f^{-1}(V)$. Consider the morphism of affines
$$
f' :
\coprod\nolimits_{i = 1, \ldots, n} U_i
\longrightarrow
V.
$$
A disjoint union of affines is affine, see
Schemes, Lemma \ref{schemes-lemma-disjoint-union-affines}.
Generic points of irreducible components of $V$
are exactly the generic points of the irreducible components of
$S$ that meet $V$. Also, $f$ is dominant if and only if $f'$ is dominant
no matter what choices of $V, n, U_i$ we make above. Thus we
have reduced the lemma to the case of a morphism of affine schemes.
The affine case is
Algebra, Lemma \ref{algebra-lemma-image-dense-generic-points}.
\end{proof}

\begin{lemma}
\label{lemma-base-change-dominant}
Let $f : X \to S$ be a quasi-compact dominant morphism of schemes.
Let $g : S' \to S$ be a morphism of schemes,  and let $f' : X' \to S'$
be the base change of $f$ by $g$. If generalizations lift along $g$,
then $f'$ is dominant.
\end{lemma}

\noindent
We will see that generalizations lift along $g$ if $g$ is open
(Lemma \ref{lemma-open-generizing})
or flat (Lemma \ref{lemma-generalizations-lift-flat}).

\begin{proof}
Observe that $f'$ is quasi-compact by Schemes, Lemma
\ref{schemes-lemma-quasi-compact-preserved-base-change}.
Let $\eta' \in S'$ be the generic point of an irreducible
component of $S'$. If generalizations lift along $g$,
then $\eta = g(\eta')$ is the generic point of an irreducible
component of $S$. By Lemma \ref{lemma-quasi-compact-dominant}
we see that $\eta$ is in the image of $f$.
Hence $\eta'$ is in the image of $f'$ by
Schemes, Lemma \ref{schemes-lemma-points-fibre-product}.
It follows that $f'$ is dominant by Lemma \ref{lemma-quasi-compact-dominant}.
\end{proof}

\noindent
Part of this holds for arbitrary (not necessarily quasi-compact)
dominant morphisms.

\begin{lemma}
\label{lemma-open-base-change-dominant}
Let $f : X \to S$ be a dominant morphism of schemes.
Let $g : S' \to S$ be an open morphism of schemes, and let $f' : X' \to S'$
be the base change of $f$ by $g$. Then $f'$ is dominant.
\end{lemma}

\begin{proof}
If $f'$ is not dominant, then $f'(X')$ is contained in some closed
subset $B$ not equal to $S'$. Then $S'-B$ is open in $S'$ and nonempty.
Since $g$ is open, $g(S'-B)$ is nonempty and open in $S$. Since $f(X)$
is dense in $S$, there is a point $x\in X$ with $f(x) \in g(S'-B)$.
Let $W = \Spec(\kappa (f(x)))$,
and write $X_W$, $S_W$, $S'_W$ for the pullbacks
of these schemes from $S$ to $W$. Then $X_W$ and $(S'-B)_W$ are nonempty,
and so $X_W \times_W (S'-B)_W$ is nonempty, as $W$ is the spectrum
of a field. This contradicts that $f'(X \times_S S')$ is contained
in $B \subset S'$. So in fact $f'$ is dominant.
\end{proof}

\begin{lemma}
\label{lemma-quasi-compact-generic-point-not-in-image}
Let $f : X \to S$ be a quasi-compact morphism of schemes.
Let $\eta \in S$ be a generic point of an irreducible
component of $S$. If $\eta \not \in f(X)$ then there
exists an open neighbourhood $V \subset S$ of $\eta$
such that $f^{-1}(V) = \emptyset$.
\end{lemma}

\begin{proof}
Let $Z \subset S$ be the scheme theoretic image of $f$.
We have to show that $\eta \not \in Z$.
This follows from
Lemma \ref{lemma-reach-points-scheme-theoretic-image}
but can also be seen as follows.
By Lemma \ref{lemma-quasi-compact-scheme-theoretic-image}
the morphism $X \to Z$ is dominant, which by
Lemma \ref{lemma-quasi-compact-dominant}
means all the generic points of all irreducible components of $Z$
are in the image of $X \to Z$. By assumption we see that
$\eta \not \in Z$ since $\eta$ would be the generic
point of some irreducible component of $Z$ if it were in $Z$.
\end{proof}

\noindent
There is another case where dominant is the same as having all
generic points of irreducible components in the image.

\begin{lemma}
\label{lemma-dominant-finite-number-irreducible-components}
Let $f : X \to S$ be a morphism of schemes.
Suppose that $X$ has finitely many irreducible components.
Then $f$ is dominant (if and) only if for every irreducible
component $Z \subset S$ the generic point of $Z$ is in the
image of $f$. If so, then $S$ has finitely many irreducible
components as well.
\end{lemma}

\begin{proof}
Assume $f$ is dominant.
Say $X = Z_1 \cup Z_2 \cup \ldots \cup Z_n$ is the decomposition
of $X$ into irreducible components. Let $\xi_i \in Z_i$ be
its generic point, so $Z_i = \overline{\{\xi_i\}}$.
Note that $f(Z_i)$ is an irreducible subset of $S$.
Hence
$$
S = \overline{f(X)} = \bigcup \overline{f(Z_i)} =
\bigcup \overline{\{f(\xi_i)\}}
$$
is a finite union of irreducible subsets whose generic
points are in the image of $f$. The lemma follows.
\end{proof}

\begin{lemma}
\label{lemma-dominant-between-integral}
Let $f : X \to Y$ be a morphism of integral schemes. The following
are equivalent
\begin{enumerate}
\item $f$ is dominant,
\item $f$ maps the generic point of $X$ to the generic point of $Y$,
\item for some nonempty affine opens $U \subset X$ and $V \subset Y$
with $f(U) \subset V$ the ring map $\mathcal{O}_Y(V) \to \mathcal{O}_X(U)$
is injective,
\item for all nonempty affine opens $U \subset X$ and $V \subset Y$
with $f(U) \subset V$ the ring map $\mathcal{O}_Y(V) \to \mathcal{O}_X(U)$
is injective,
\item for some $x \in X$ with image $y = f(x) \in Y$ the local ring
map $\mathcal{O}_{Y, y} \to \mathcal{O}_{X, x}$ is injective, and
\item for all $x \in X$ with image $y = f(x) \in Y$ the local ring
map $\mathcal{O}_{Y, y} \to \mathcal{O}_{X, x}$ is injective.
\end{enumerate}
\end{lemma}

\begin{proof}
The equivalence of (1) and (2) follows from
Lemma \ref{lemma-dominant-finite-number-irreducible-components}.
Let $U \subset X$ and $V \subset Y$ be nonempty affine opens with
$f(U) \subset V$. Recall that the rings $A = \mathcal{O}_X(U)$
and $B = \mathcal{O}_Y(V)$ are integral domains.
The morphism $f|_U : U \to V$ corresponds to a ring map
$\varphi : B \to A$. The generic points of $X$ and $Y$ correspond
to the prime ideals $(0) \subset A$ and $(0) \subset B$. Thus
(2) is equivalent to the condition $(0) = \varphi^{-1}((0))$,
i.e., to the condition that $\varphi$ is injective.
In this way we see that (2), (3), and (4) are equivalent.
Similarly, given $x$ and $y$ as in (5) the local rings
$\mathcal{O}_{X, x}$ and $\mathcal{O}_{Y, y}$ are domains
and the prime ideals $(0) \subset \mathcal{O}_{X, x}$
and $(0) \subset \mathcal{O}_{Y, y}$ correspond to the
generic points of $X$ and $Y$ (via the identification of
the spectrum of the local ring at $x$
with the set of points specializing to $x$, see
Schemes, Lemma \ref{schemes-lemma-specialize-points}).
Thus we can argue in the exact same manner as above
to see that (2), (5), and (6) are equivalent.
\end{proof}







